\documentclass{article}
\usepackage{graphicx} % Required for inserting images

\title{qmpbs}
\author{jpbruneton }
\date{July 2026}

\begin{document}

\maketitle

% Requires:
% \usepackage{enumitem}
% \usepackage{xcolor}
%
% \newcommand{\precise}{\textcolor{blue!70!black}{\textsf{[P]}}}
% \newcommand{\formalism}{\textcolor{orange!80!black}{\textsf{[F]}}}

\section{Open problems in theoretical quantum theory}

The label \precise{} denotes a problem that can be formulated as a
mathematical conjecture, decision problem, or classification problem.
The label \formalism{} denotes a sharply identifiable gap in the
formalism, even when no unique conjectural statement has yet emerged.

\subsection{Mathematical quantum mechanics and spectral theory}

\begin{enumerate}[label=\textbf{M\arabic*.},leftmargin=*]

\item \precise{} \textbf{Ionization conjecture for atoms.}
For the nonrelativistic Coulomb Hamiltonian of a nucleus of charge \(Z\),
prove that the maximal number \(N_{\mathrm c}(Z)\) of electrons which can
be bound satisfies
\[
    N_{\mathrm c}(Z)\leq Z+C
\]
with a universal constant \(C\), independent of \(Z\).  The corresponding
statement is known for several mean-field models, but not for the full
many-electron Schrödinger Hamiltonian.

\item \precise{} \textbf{Delocalization in the three-dimensional Anderson model.}
Prove the existence of an absolutely continuous spectral regime and of
extended states for the standard random Schrödinger operator on
\(\mathbb Z^{3}\) at sufficiently weak disorder, and establish rigorously
the localization--delocalization transition.

\item \precise{} \textbf{Mobility edges.}
For a natural random Schrödinger operator, prove the existence and
determine the location or regularity of an energy separating localized
from delocalized states.

\item \precise{} \textbf{Many-body localization.}
Prove or disprove the existence of a stable many-body-localized phase for
a physically natural interacting lattice Hamiltonian in the
thermodynamic limit, including stability under generic sufficiently
small local perturbations.

\item \precise{} \textbf{The Haldane gap for the Heisenberg chain.}
Prove that the translation-invariant antiferromagnetic spin-\(1\)
Heisenberg Hamiltonian
\[
    H=\sum_{j\in\mathbb Z}\vec S_j\cdot\vec S_{j+1}
\]
has a nonzero spectral gap above its ground state in the thermodynamic
limit.  More generally, settle the integer-spin part of Haldane's
conjecture for the standard Heisenberg chains.

\item \precise{} \textbf{Gap stability beyond the presently controlled classes.}
Find general, checkable hypotheses under which the bulk spectral gap of
an interacting infinite quantum system is stable under local
perturbations, including systems with long-range interactions,
degenerate ground spaces, or gapless boundary modes.

\item \precise{} \textbf{Spectral gaps of PEPS parent Hamiltonians.}
Develop effective criteria deciding whether the parent Hamiltonian of a
specified family of projected entangled-pair states is uniformly
gapped in the thermodynamic limit.

\item \precise{} \textbf{Quantum unique ergodicity for chaotic systems.}
Determine natural conditions under which every high-energy sequence of
eigenfunctions of a classically ergodic Hamiltonian becomes
equidistributed, rather than merely a density-one subsequence.

\item \precise{} \textbf{Bohigas--Giannoni--Schmit conjecture.}
Prove that the local spectral statistics of a generic quantum system
whose classical limit is chaotic converge to the appropriate random
matrix ensemble.

\item \precise{} \textbf{Scars in the semiclassical limit.}
Characterize the possible semiclassical measures supported partly on
unstable periodic orbits and determine when positive-density or
exceptional scarred subsequences can exist.

\item \formalism{} \textbf{Infinite-dimensional quantum information.}
Develop a unified theory of states, channels, entropies and resource
conversion on infinite-dimensional Hilbert spaces which handles domain
questions, unbounded observables and energy constraints without
reducing every result to an ad hoc finite-dimensional truncation.

\end{enumerate}


\subsection{Quantum many-body theory}

\begin{enumerate}[label=\textbf{B\arabic*.},leftmargin=*]

\item \precise{} \textbf{Area law in two spatial dimensions.}
Prove that the ground state of every sufficiently local,
finite-dimensional, gapped two-dimensional Hamiltonian satisfies
\[
    S(\rho_A)\leq C\,|\partial A|
\]
up to a precisely controlled correction, without assuming frustration
freeness or an adiabatic connection to a trivial phase.

\item \precise{} \textbf{Efficient tensor-network approximation in two dimensions.}
Determine whether every gapped two-dimensional ground state admits a
PEPS or another efficiently describable tensor-network approximation
with polynomial bond dimension for local observables.

\item \precise{} \textbf{Efficient contraction of physically restricted PEPS.}
Identify the largest natural class of gapped or injective PEPS for which
norms and local expectation values can be approximated in polynomial
time.

\item \precise{} \textbf{Eigenstate thermalization hypothesis.}
Give checkable hypotheses on a nonintegrable local Hamiltonian implying
the diagonal and off-diagonal ETH estimates for individual
finite-energy-density eigenstates.

\item \precise{} \textbf{Unitary derivation of thermalization.}
Prove thermalization, with quantitative time scales and finite-size
errors, for a nonintegrable translation-invariant many-body Hamiltonian
and a broad class of nonequilibrium initial states.

\item \precise{} \textbf{Derivation of quantum hydrodynamics.}
Starting from local unitary many-body dynamics, derive diffusion,
ballistic transport or nonlinear fluctuating hydrodynamics, including
explicit error bounds in an appropriate scaling limit.

\item \formalism{} \textbf{Classification of failures of thermalization.}
Construct a common framework encompassing integrability, many-body
localization, Hilbert-space fragmentation, quantum scars and conserved
pseudolocal charges, and determine which mechanisms are stable phases
rather than fine-tuned exceptions.

\item \precise{} \textbf{Classification of gapped phases in \(d\geq2\).}
Classify gapped local Hamiltonians up to gap-preserving local
deformations, including intrinsic topological order, symmetry-enriched
phases, chiral phases and anomalous boundaries.

\item \precise{} \textbf{Completeness of topological invariants.}
Determine whether known invariants---anyon data, fusion and braiding
categories, chiral central charge, symmetry fractionalization and
higher-form data---completely classify the corresponding gapped phases.

\item \precise{} \textbf{Two-dimensional Hubbard model.}
Determine rigorously the zero-temperature phase diagram of the
repulsive Hubbard Hamiltonian on the square lattice, in particular the
existence and relation of antiferromagnetism, pseudogap behaviour and
superconducting order away from half filling.

\item \precise{} \textbf{Intrinsic sign problem.}
Characterize which local Hamiltonians can be transformed, by an
efficiently describable local or quasilocal change of basis, into a
stoquastic Hamiltonian.

\item \precise{} \textbf{Optimal Lieb--Robinson bounds for long-range systems.}
Determine the sharp effective causal region and information-propagation
speed for interactions decaying as \(r^{-\alpha}\), in all regimes of
\(\alpha\) and spatial dimension.

\end{enumerate}


\subsection{Mathematical quantum field theory}

\begin{enumerate}[label=\textbf{QF\arabic*.},leftmargin=*]

\item \precise{} \textbf{Yang--Mills existence and mass gap.}
For every compact simple gauge group \(G\), construct a nontrivial
quantum Yang--Mills theory on \(\mathbb R^{4}\) satisfying suitable
axioms and prove that its Hamiltonian has a mass gap
\[
    \Delta>0 .
\]

\item \precise{} \textbf{Confinement in four-dimensional Yang--Mills theory.}
Prove an appropriate confinement criterion, such as an area law for
large Wilson loops, and connect it rigorously with the physical
spectrum and colour-singlet structure.

\item \precise{} \textbf{Construction of four-dimensional QED.}
Construct continuum quantum electrodynamics nonperturbatively and prove
that it satisfies a precise set of relativistic QFT axioms, while
controlling charge renormalization, infrared sectors and the possible
Landau-pole problem.

\item \precise{} \textbf{Nonperturbative construction of the Standard Model.}
Give a mathematically complete continuum definition of the chiral gauge
theory underlying the Standard Model and establish its consistency
beyond formal perturbation theory.

\item \precise{} \textbf{Chiral gauge theories on the lattice.}
Construct a general nonperturbative lattice regularization of
anomaly-free chiral gauge theories whose continuum limit has the
required locality, gauge symmetry and fermion content.

\item \precise{} \textbf{Asymptotic completeness in interacting relativistic QFT.}
For a nontrivial \(3+1\)-dimensional interacting QFT, prove that its
physical Hilbert space is exhausted by the vacuum, stable particles and
their scattering states, with a suitable treatment of massless
particles and infraparticles.

\item \formalism{} \textbf{Infrared sectors and infraparticles.}
Develop a complete scattering formalism for theories such as QED in
which charged states are accompanied by infinitely many soft quanta and
do not belong to ordinary sharp-mass Wigner particle sectors.

\item \formalism{} \textbf{Entanglement in continuum QFT.}
Give a regulator-independent and operationally meaningful general
formalism for entanglement and resource conversion when local
observable algebras are type III and density matrices for sharp spatial
regions do not exist.

\item \formalism{} \textbf{Subsystems in gauge theories.}
Find a canonical treatment of spatial subsystems, edge modes, centres
of local algebras and gauge constraints, clarifying which definitions
of entanglement correspond to distinct operational tasks.

\item \precise{} \textbf{Interacting QFT on curved spacetime.}
Construct nonperturbative interacting quantum fields on broad classes
of curved globally hyperbolic spacetimes and characterize physically
distinguished states without relying on a preferred vacuum.

\end{enumerate}


\subsection{Entanglement theory}

\begin{enumerate}[label=\textbf{E\arabic*.},leftmargin=*]

\item \precise{} \textbf{NPT bound entanglement.}
Determine whether every bipartite state with negative partial transpose
is distillable, or construct an NPT state whose distillable
entanglement is zero.

\item \precise{} \textbf{Distillability of NPT Werner states.}
For the unresolved range of NPT Werner states, decide whether there
exist an integer \(n\) and a Schmidt-rank-two vector
\(\lvert\psi\rangle\) such that
\[
 \langle\psi|(\rho_W^{T_B})^{\otimes n}|\psi\rangle<0 .
\]

\item \precise{} \textbf{Effective separability criteria.}
Find complete and computationally useful separability criteria for
nontrivial families beyond \(2\times2\) and \(2\times3\), while
identifying the maximal classes for which separability is tractable.

\item \precise{} \textbf{Exact characterization of LOCC.}
Characterize the closure and boundary of the set of transformations
implementable by local operations and classical communication, and
give an effective membership criterion.

\item \precise{} \textbf{Asymptotic mixed-state entanglement conversion.}
Determine the optimal asymptotic rate
\[
    R(\rho\to\sigma)
\]
for general bipartite mixed states under LOCC, including the possible
role of catalysts and sublinear auxiliary resources.

\item \precise{} \textbf{Irreversibility of mixed-state entanglement.}
Characterize exactly when entanglement cost and distillable
entanglement coincide:
\[
    E_{\mathrm C}(\rho)=E_{\mathrm D}(\rho).
\]

\item \precise{} \textbf{Minimal reversible entanglement-generating set.}
Determine whether there exists a finite minimal reversible
entanglement-generating set for asymptotic multipartite LOCC
transformations.

\item \precise{} \textbf{Multipartite entanglement classification.}
Find a scalable classification of multipartite mixed states under LOCC
or SLOCC, including a useful organization of the infinitely many
inequivalent families that appear as the number of parties grows.

\item \precise{} \textbf{Computability of regularized entanglement measures.}
Determine which quantities such as
\[
 E_{\mathrm C},\qquad E_{\mathrm D},\qquad
 E_R^\infty,\qquad E_P^\infty
\]
are computable or approximable with certified error for general
finite-dimensional states.

\item \precise{} \textbf{Additivity of entanglement of purification.}
Determine whether
\[
    E_P(\rho\otimes\sigma)
      =E_P(\rho)+E_P(\sigma)
\]
holds, or construct an explicit counterexample and determine the
regularized quantity \(E_P^\infty\).

\item \precise{} \textbf{Entanglement catalysis.}
Give necessary and sufficient conditions for catalytic state conversion
when the states and catalysts may be mixed, multipartite or correlated,
and determine the resources required of the catalyst.

\item \precise{} \textbf{Quantum marginal problem.}
Give a complete, effectively usable characterization of collections
\(\{\rho_{A_i}\}\) that arise as compatible reduced states of one global
state, beyond the cases already covered by representation-theoretic
solutions.

\item \precise{} \textbf{Quantum entropy cone.}
Characterize the closure of the set
\[
 \bigl\{(S(\rho_I))_{\varnothing\neq I\subseteq[n]}\bigr\}
\]
for \(n\geq4\), including whether there exist new independent linear
entropy inequalities beyond the presently known universal ones.

\item \precise{} \textbf{Absolutely maximally entangled states.}
Classify the pairs \((n,d)\) for which an
\(\operatorname{AME}(n,d)\) state exists, and give constructive
existence or nonexistence proofs.

\item \formalism{} \textbf{Entanglement of identical particles.}
Establish a common operational framework reconciling mode
entanglement, particle entanglement, superselection rules and
exchange correlations for bosons and fermions.

\end{enumerate}


\subsection{Quantum measurements, nonlocality and contextuality}

\begin{enumerate}[label=\textbf{N\arabic*.},leftmargin=*]

\item \precise{} \textbf{Existence of SIC-POVMs.}
Prove or disprove that for every \(d\geq2\) there exist \(d^{2}\) unit
vectors \(\{\lvert\psi_j\rangle\}\subset\mathbb C^d\) satisfying
\[
    |\langle\psi_j|\psi_k\rangle|^2=\frac{1}{d+1},
    \qquad j\neq k .
\]

\item \precise{} \textbf{Mutually unbiased bases in dimension six.}
Determine the maximal number \(M(6)\) of mutually unbiased orthonormal
bases in \(\mathbb C^6\).  In particular, decide whether
\[
    M(6)=3
\]
or whether a set of four or more such bases exists.

\item \precise{} \textbf{MUBs in non-prime-power dimensions.}
Determine the maximal number \(M(d)\) of mutually unbiased bases when
\(d\) is not a prime power, and decide whether complete sets of
\(d+1\) bases can exist in any such dimension.

\item \precise{} \textbf{Minimal separation of quantum correlation models.}
Find the smallest Bell scenario in which finite-dimensional,
approximately finite-dimensional and commuting-operator quantum
correlation sets differ.

\item \precise{} \textbf{Dimension required for nonlocal correlations.}
Given a Bell correlation and accuracy \(\varepsilon\), bound or
determine the minimal local Hilbert-space dimension needed to realize
it within \(\varepsilon\).

\item \precise{} \textbf{Restricted decidability of nonlocal games.}
Although the general problem is undecidable, classify the natural
families of nonlocal games for which the quantum value is computable,
semidecidable or efficiently approximable.

\item \precise{} \textbf{Classification of self-testing correlations.}
Characterize which states and measurements can be uniquely certified,
up to local isometries, from their correlations alone.

\item \precise{} \textbf{Optimal robustness of self-testing.}
For standard self-tests, determine the sharp relation between the
deviation of a Bell value from its optimum and the distance from the
ideal state and measurements.

\item \precise{} \textbf{Quantum maxima of Bell and contextuality inequalities.}
Develop exact methods for calculating the quantum maximum and the
minimal realizing dimension of a general inequality, beyond
convergent but potentially nonterminating semidefinite hierarchies.

\item \formalism{} \textbf{Contextuality as a computational resource.}
Find necessary and sufficient conditions under which contextuality
provides an actual quantum computational advantage, rather than merely
being present in a computation.

\end{enumerate}


\subsection{Quantum channels and quantum Shannon theory}

\begin{enumerate}[label=\textbf{C\arabic*.},leftmargin=*]

\item \precise{} \textbf{Computable classes of quantum capacities.}
Because no general single-letter formula can exist in full generality,
identify maximal natural classes of channels for which the quantum
capacity
\[
    Q(\mathcal N)
      =\lim_{n\to\infty}\frac1n
       \max_{\rho} I_{\mathrm c}(\rho,\mathcal N^{\otimes n})
\]
is computable or admits a finite-letter expression.

\item \precise{} \textbf{Quantum capacity of the thermal attenuator.}
Determine the energy-constrained quantum and private capacities of a
general bosonic thermal-loss channel outside the degradable and
antidegradable regimes.

\item \precise{} \textbf{Two-way capacities of Gaussian channels.}
Determine the exact entanglement-distribution and secret-key
capacities assisted by unlimited two-way classical communication for
thermal attenuators, amplifiers and additive-noise channels.

\item \precise{} \textbf{Strong converse for quantum capacity.}
Determine for which channels communication above \(Q(\mathcal N)\)
forces the fidelity to converge to zero, and settle the general
finite-dimensional case or provide a counterexample.

\item \precise{} \textbf{Private versus quantum capacity.}
Characterize channels for which
\[
    P(\mathcal N)>Q(\mathcal N),
\]
determine the maximal possible separation, and decide when either
capacity is additive.

\item \precise{} \textbf{Zero-error quantum capacities.}
Characterize zero-error classical and quantum capacity, their
superactivation, and the channel classes for which the required
regularization can be evaluated effectively.

\item \precise{} \textbf{Quantum broadcast-channel capacity region.}
Determine the full classical--quantum--entanglement capacity region of
a general quantum broadcast channel.

\item \precise{} \textbf{Quantum interference-channel capacity region.}
Find matching inner and outer bounds, or an exact characterization, for
general quantum interference channels.

\item \precise{} \textbf{Capacities of channels with memory.}
Develop computable capacity formulas for broad non-Markovian channel
families with correlated noise, including a precise operational
classification of forgetful and nonforgetful memory.

\item \formalism{} \textbf{Quantum Shannon theory in infinite dimensions.}
Extend coding theorems, strong converses and resource inequalities to
bosonic and other infinite-dimensional systems under physically
natural energy constraints, with controlled compactness and continuity
assumptions.

\end{enumerate}


\subsection{Quantum algorithms and complexity}

\begin{enumerate}[label=\textbf{A\arabic*.},leftmargin=*]

\item \precise{} \textbf{\(\mathsf{BQP}\) versus classical computation.}
Prove an unrelativized separation such as
\[
    \mathsf{BPP}\neq\mathsf{BQP},
\]
or otherwise characterize exactly which classical complexity class
captures efficient quantum computation.

\item \precise{} \textbf{\(\mathsf{NP}\) versus \(\mathsf{BQP}\).}
Determine whether
\[
    \mathsf{NP}\subseteq\mathsf{BQP},
\]
and in particular whether quantum computers can solve
\(\mathsf{NP}\)-complete problems in polynomial time.

\item \precise{} \textbf{\(\mathsf{QMA}\) versus \(\mathsf{QCMA}\).}
Determine whether quantum witnesses are strictly more powerful than
classical witnesses for polynomial-time quantum verification:
\[
    \mathsf{QMA}\stackrel{?}{=}\mathsf{QCMA}.
\]

\item \precise{} \textbf{Perfect completeness for QMA.}
Determine whether
\[
    \mathsf{QMA}_1=\mathsf{QMA}.
\]

\item \precise{} \textbf{Multiple unentangled quantum proofs.}
Determine whether
\[
    \mathsf{QMA}(2)=\mathsf{QMA},
\]
or whether two unentangled witnesses strictly increase verification
power.

\item \precise{} \textbf{Quantum PCP conjecture.}
Prove or disprove that approximating the ground-state energy of a local
Hamiltonian to constant extensive precision is
\(\mathsf{QMA}\)-hard.

\item \precise{} \textbf{Good quantum locally testable codes.}
Construct, or rule out, quantum locally testable codes having
simultaneously constant rate, linear distance and constant soundness
with bounded-weight checks.

\item \precise{} \textbf{Commuting local Hamiltonian problem.}
Determine the complexity of the commuting local Hamiltonian problem
for general fixed local dimension and interaction geometry, and whether
all relevant cases admit succinct classical witnesses.

\item \precise{} \textbf{Stoquastic Hamiltonian complexity.}
Determine the relation between \(\mathsf{StoqMA}\), \(\mathsf{MA}\) and
related classical classes; in particular, settle whether general
error amplification is possible for \(\mathsf{StoqMA}\).

\item \precise{} \textbf{Quantum search-to-decision reductions.}
Determine whether a general accepting quantum witness can be prepared
efficiently using only polynomially many queries to a decision oracle
for the corresponding \(\mathsf{QMA}\) problem.

\item \precise{} \textbf{Aaronson--Ambainis conjecture.}
Prove that every bounded low-degree polynomial on the Boolean cube has
an influential variable of polynomially bounded influence, with the
conjectured consequences for classical simulation of quantum query
algorithms on most inputs.

\item \precise{} \textbf{Average-case hardness of BosonSampling.}
Complete the worst-to-average-case reduction and anticoncentration
arguments needed to establish approximate classical sampling hardness
under standard complexity assumptions.

\item \precise{} \textbf{Average-case hardness of random-circuit sampling.}
Prove approximate sampling hardness for natural random-circuit
ensembles using assumptions no stronger than standard worst-case
complexity conjectures.

\item \precise{} \textbf{Graph isomorphism and quantum algorithms.}
Determine whether graph isomorphism belongs to \(\mathsf{BQP}\), or
prove meaningful quantum lower bounds ruling out broad classes of
quantum algorithms.

\item \precise{} \textbf{Lattice problems and quantum computation.}
Determine the exact quantum complexity of standard lattice problems
such as approximate shortest vector and closest vector problems in the
parameter regimes relevant to classical and post-quantum cryptography.

\item \formalism{} \textbf{A criterion for genuine quantum speedup.}
Develop structural conditions on a computational problem that predict
whether it admits more than a polynomial quantum advantage, while
excluding speedups caused only by input, output or data-access models.

\end{enumerate}


\subsection{Computability and undecidability}

\begin{enumerate}[label=\textbf{U\arabic*.},leftmargin=*]

\item \precise{} \textbf{Boundary between decidable and undecidable gap problems.}
Classify the restrictions on dimension, translation invariance, local
dimension, interaction range and frustration under which the
thermodynamic spectral-gap problem becomes decidable.

\item \precise{} \textbf{Computability of phase diagrams.}
Determine which families of local Hamiltonians have computable phase
boundaries and order parameters, given that unrestricted phase
questions can encode undecidable problems.

\item \precise{} \textbf{Computability of thermalization.}
Classify the local Hamiltonian families for which long-time averages,
thermalization and equilibration are decidable or effectively
approximable.

\item \precise{} \textbf{Computability of channel capacities.}
Classify quantum channels according to whether their classical,
private, quantum and zero-error capacities are computable,
semicomputable or noncomputable.

\item \formalism{} \textbf{Computable analysis for quantum mechanics.}
Develop a physically appropriate computability framework for
unbounded operators, continuous spectra, generalized eigenstates and
infinite tensor products, including a clear account of how input data
are specified with finite information.

\end{enumerate}


\subsection{Open quantum systems and quantum thermodynamics}

\begin{enumerate}[label=\textbf{O\arabic*.},leftmargin=*]

\item \precise{} \textbf{Unbounded Lindblad generators.}
Give necessary and sufficient conditions characterizing the
generators of completely positive, trace-preserving semigroups on
infinite-dimensional systems when the Hamiltonian and jump operators
are unbounded.

\item \precise{} \textbf{Controlled derivation of master equations.}
Derive Markovian or non-Markovian reduced dynamics from a microscopic
system--environment Hamiltonian with explicit error bounds valid on
specified time scales.

\item \formalism{} \textbf{Strong-coupling open-system dynamics.}
Construct a general reduced-dynamics formalism valid beyond weak
coupling, secular approximations and initially factorized states,
while preserving positivity and thermodynamic consistency.

\item \precise{} \textbf{Initial system--environment correlations.}
Characterize exactly which reduced dynamical maps are physically
admissible for specified families of initially correlated
system--environment states.

\item \formalism{} \textbf{Quantum non-Markovianity.}
Find an operational definition of memory that unifies or clearly
relates divisibility, information backflow, process tensors,
multi-time correlations and simulation cost.

\item \precise{} \textbf{Thermal operations with coherence.}
Give necessary and sufficient conditions for state conversion under
thermal operations when the states contain coherence between energy
eigenspaces, with and without correlated catalysts.

\item \precise{} \textbf{Autonomous quantum machines.}
Determine the exact performance limits of clocks, batteries,
refrigerators and engines implemented by time-independent finite
Hamiltonians, including the cost of control and degradation of the
auxiliary systems.

\item \formalism{} \textbf{Work and heat at strong coupling.}
Construct an operational decomposition of energy changes into work and
heat that is consistent for strongly coupled, correlated and
non-Markovian quantum systems.

\end{enumerate}


\subsection{Foundations and missing formalism}

\begin{enumerate}[label=\textbf{F\arabic*.},leftmargin=*]

\item \formalism{} \textbf{Measurement problem.}
Give a mathematically precise account of individual definite outcomes
that is compatible with the successful unitary dynamics of quantum
theory and does not introduce an undefined boundary between system,
apparatus and observer.

\item \formalism{} \textbf{Born rule from explicit assumptions.}
Derive the Born probability rule from a minimal set of clearly stated
physical or decision-theoretic assumptions without presupposing
probabilistic, noncontextual or branch-weight structure equivalent to
the desired conclusion.

\item \formalism{} \textbf{Operational reconstruction of quantum theory.}
Find a minimal and conceptually compelling set of operational
principles which uniquely produces complex Hilbert spaces, tensor
products, completely positive maps and the Born rule, while explaining
why real and quaternionic alternatives are excluded.

\item \formalism{} \textbf{Wigner's-friend scenarios.}
Provide a consistent multi-agent probability calculus specifying when
different agents may combine statements about outcomes, including
cases where laboratories and observers are treated quantum
mechanically.

\item \formalism{} \textbf{Emergence of classical outcomes.}
Explain quantitatively how a preferred set of stable classical records
and effectively Boolean events emerges from unitary quantum dynamics,
including the role of decoherence, redundancy and coarse graining.

\item \formalism{} \textbf{Classical limit.}
State and prove sufficiently general conditions under which quantum
dynamics, probabilities and histories converge to an appropriate
classical theory, including chaotic and macroscopic systems.

\item \formalism{} \textbf{Definition of a quantum subsystem.}
Characterize when a tensor-product factorization
\[
    \mathcal H\simeq\mathcal H_A\otimes\mathcal H_B
\]
has physical significance, how it is selected by observable algebras
and interactions, and when inequivalent factorizations describe
different notions of locality and entanglement.

\item \formalism{} \textbf{Quantum reference frames.}
Develop a fully compositional formalism for changing between quantum
reference frames, including imperfect, finite, entangled and
relativistic frames and systems subject to gauge constraints.

\item \formalism{} \textbf{Physical boundary of indefinite causal order.}
Characterize which mathematically valid process matrices can arise
from ordinary quantum systems embedded in spacetime and which require
a genuine extension of standard quantum theory.

\item \formalism{} \textbf{Quantum causal models.}
Find a general quantum analogue of classical causal inference that
combines interventions, latent quantum systems, indefinite order and
noncommuting observables while retaining identifiable causal
statements.

\item \formalism{} \textbf{Time observables and quantum clocks.}
Give a unified treatment of time as an external parameter, a POVM, a
dynamical clock and a relational observable, including the limitations
caused by finite clock energy and backreaction.

\end{enumerate}
\end{document}
